\documentclass[11pt,a4paper]{article}

\def\courseTitle{Industrial Security (Bachelor)}
\def\labNumber{03}
\def\labTitle{Network Architecture -- Purdue, Zones, and Conduits}
\def\labSection{03}
\def\labTime{75--90 min}

% =====================================================================
% Shared preamble for lab sheets.
%
% Caller must \documentclass{article} first, then define:
%   \def\courseTitle{...}
%   \def\labNumber{NN}
%   \def\labTitle{...}
%   \def\labTime{e.g. "30--45 min"}
%   \def\labSection{e.g. "02"}
% Then % =====================================================================
% Shared preamble for lab sheets.
%
% Caller must \documentclass{article} first, then define:
%   \def\courseTitle{...}
%   \def\labNumber{NN}
%   \def\labTitle{...}
%   \def\labTime{e.g. "30--45 min"}
%   \def\labSection{e.g. "02"}
% Then % =====================================================================
% Shared preamble for lab sheets.
%
% Caller must \documentclass{article} first, then define:
%   \def\courseTitle{...}
%   \def\labNumber{NN}
%   \def\labTitle{...}
%   \def\labTime{e.g. "30--45 min"}
%   \def\labSection{e.g. "02"}
% Then \input{../../template/lab-preamble.tex}
% =====================================================================

\usepackage[utf8]{inputenc}
\usepackage[T1]{fontenc}
\usepackage[english]{babel}
\usepackage[a4paper, margin=2cm]{geometry}
\usepackage{graphicx}
\usepackage{listings}
\usepackage{xcolor}
\usepackage{colortbl}
\usepackage{tikz}
\usepackage{amsmath}
\usepackage{amssymb}
\usepackage{booktabs}
\usepackage{enumitem}
\usepackage{titlesec}
\usepackage{fancyhdr}
\usepackage{hyperref}
\usepackage{tcolorbox}
\tcbuselibrary{skins, breakable}
\usepackage{fontawesome5}

\input{../../../template/tikz-styles.tex}

\providecommand{\courseAuthor}{Industrial Security Lecture Series}
\providecommand{\courseInstitute}{Department of Computer Science}
\providecommand{\companyLogo}{logo}

\definecolor{slideAccent}{HTML}{00205B}
\definecolor{slideSecondary}{HTML}{00A0DC}

% --- Corporate branding overrides ------------------------------------
\IfFileExists{../../../branding.tex}{\input{../../../branding.tex}}{}

\colorlet{labAccent}{slideAccent}
\colorlet{labSec}{slideSecondary}
\definecolor{labCheck}{HTML}{2E7D32}
\definecolor{labWarn}{HTML}{B23A48}

\lstset{
  backgroundcolor=\color{black!4},
  basicstyle=\ttfamily\footnotesize,
  breaklines=true,
  frame=leftline,
  framerule=2pt,
  rulecolor=\color{labAccent},
  numbers=left,
  numberstyle=\tiny\color{black!50},
  showstringspaces=false,
  tabsize=2
}

\setlength{\tabcolsep}{4pt}

% Let TeX add a little extra inter-word stretch as a last resort before
% it reports an Overfull \hbox. Only kicks in on lines that would
% otherwise overflow (long \texttt paths, URLs, CIDRs); never loosens a
% line that already fits. Keeps prose/itemize within the margin.
\setlength{\emergencystretch}{3em}

\pagestyle{fancy}
\fancyhf{}
\renewcommand{\headrulewidth}{0.4pt}
\lhead{\small \courseTitle{} -- Lab \labNumber}
\rhead{\small Maps to Section \labSection}
\cfoot{\thepage}

\newtcolorbox{stepbox}[1]{
  enhanced, breakable, arc=2pt, boxrule=0pt,
  colback=labAccent!7, colframe=labAccent, leftrule=3pt,
  fonttitle=\bfseries\small, coltitle=white,
  colbacktitle=labAccent,
  title={\faTools\ Step #1}
}

\newtcolorbox{warnbox}{
  enhanced, breakable, arc=2pt, boxrule=0pt,
  colback=labWarn!7, colframe=labWarn, leftrule=4pt,
  fonttitle=\bfseries\small, coltitle=white, colbacktitle=labWarn,
  title={\faExclamationTriangle\ Caution}
}

\newtcolorbox{notebox}{
  enhanced, breakable, arc=2pt, boxrule=0pt,
  colback=labAccent!4, colframe=labAccent, leftrule=2pt,
  fonttitle=\bfseries\small, coltitle=white, colbacktitle=labAccent,
  title={\faInfoCircle\ Note}
}

\newtcolorbox{goalbox}{
  enhanced, breakable, arc=2pt, boxrule=0pt,
  colback=labCheck!7, colframe=labCheck, leftrule=3pt,
  fonttitle=\bfseries\small, coltitle=white, colbacktitle=labCheck,
  title={\faBullseye\ Goal}
}

\newtcolorbox{deliverbox}{
  enhanced, breakable, arc=2pt, boxrule=0pt,
  colback=labSec!10, colframe=labSec, leftrule=3pt,
  fonttitle=\bfseries\small, coltitle=white, colbacktitle=labSec,
  title={\faClipboardList\ Deliverable}
}

% --- Solution-sheet boxes (used only by lab-solutions.tex) ----------
\newtcolorbox{solutionbox}[1]{
  enhanced, breakable, arc=2pt, boxrule=0pt,
  colback=labCheck!7, colframe=labCheck, leftrule=3pt,
  fonttitle=\bfseries\small, coltitle=white, colbacktitle=labCheck,
  title={\faCheckCircle\ #1}
}

\newtcolorbox{rubricbox}{
  enhanced, breakable, arc=2pt, boxrule=0pt,
  colback=labSec!7, colframe=labSec, leftrule=3pt,
  fonttitle=\bfseries\small, coltitle=white, colbacktitle=labSec,
  title={\faBalanceScale\ Grading rubric}
}

\titleformat{\section}{\large\bfseries\color{labAccent}}{\thesection}{0.5em}{}

\newcommand{\labheader}{%
  \noindent
  \begin{tikzpicture}
    \fill[labAccent] (0,0) rectangle (\textwidth, -1.6cm);
    \node[anchor=north west, text=white, font=\bfseries\large] at (0.6cm, -0.4cm) {Lab \labNumber{} -- \labTitle};
    \node[anchor=south west, text=white, font=\small] at (0.6cm, -1.3cm) {\courseTitle{} \quad $\bullet$ \quad Maps to Section \labSection{} \quad $\bullet$ \quad \labTime};
    \node[anchor=north east, fill=labSec, text=white, font=\bfseries, inner sep=4pt, rounded corners=2pt]
      at ([xshift=-0.4cm, yshift=-0.4cm]\textwidth,0) {LAB \labNumber};
  \end{tikzpicture}
  \vspace{4mm}
}

% =====================================================================

\usepackage[utf8]{inputenc}
\usepackage[T1]{fontenc}
\usepackage[english]{babel}
\usepackage[a4paper, margin=2cm]{geometry}
\usepackage{graphicx}
\usepackage{listings}
\usepackage{xcolor}
\usepackage{colortbl}
\usepackage{tikz}
\usepackage{amsmath}
\usepackage{amssymb}
\usepackage{booktabs}
\usepackage{enumitem}
\usepackage{titlesec}
\usepackage{fancyhdr}
\usepackage{hyperref}
\usepackage{tcolorbox}
\tcbuselibrary{skins, breakable}
\usepackage{fontawesome5}

% =====================================================================
% Shared TikZ styles for the Industrial Security lecture series.
% Loaded by every slide deck and exercise sheet that uses TikZ.
% =====================================================================

\usetikzlibrary{
  arrows.meta,
  shapes,
  shapes.geometric,
  shapes.symbols,
  shapes.multipart,
  positioning,
  fit,
  calc,
  backgrounds,
  decorations.pathmorphing,
  decorations.pathreplacing,
  decorations.text,
  patterns,
  matrix,
  shadows
}

% --- colour palette --------------------------------------------------
\definecolor{isOT}{HTML}{1F4E79}     % OT (deep blue)
\definecolor{isIT}{HTML}{6F2DA8}     % IT (purple)
\definecolor{isSafety}{HTML}{C00000} % safety red
\definecolor{isNet}{HTML}{2E7D32}    % network green
\definecolor{isAttack}{HTML}{B23A48} % attack red
\definecolor{isAccent}{HTML}{E08E0B} % amber
\definecolor{isMute}{HTML}{6B6B6B}   % grey
\definecolor{isLight}{HTML}{EDF2F8}  % light background
\definecolor{isPLC}{HTML}{2C5282}
\definecolor{isHMI}{HTML}{2F855A}
\definecolor{isSCADA}{HTML}{6B46C1}

% --- generic node styles --------------------------------------------
\tikzset{
  isnode/.style={
    draw, rounded corners=2pt, minimum height=8mm, minimum width=18mm,
    font=\small, align=center, thick
  },
  isbox/.style={isnode, fill=isLight},
  plc/.style={isnode, fill=isPLC!15, draw=isPLC, thick},
  hmi/.style={isnode, fill=isHMI!15, draw=isHMI, thick},
  scada/.style={isnode, fill=isSCADA!15, draw=isSCADA, thick},
  sensor/.style={isnode, fill=isNet!10, draw=isNet, minimum height=6mm, minimum width=14mm, font=\scriptsize},
  actuator/.style={isnode, fill=isAccent!15, draw=isAccent, minimum height=6mm, minimum width=14mm, font=\scriptsize},
  firewall/.style={isnode, fill=isAttack!15, draw=isAttack, thick},
  server/.style={isnode, fill=isMute!15, draw=isMute!80, thick},
  switch/.style={isnode, fill=isLight, draw=black!60, minimum height=6mm, minimum width=14mm, font=\scriptsize},
  workstation/.style={isnode, fill=white, draw=isOT, thick},
  attacker/.style={isnode, fill=isAttack!20, draw=isAttack, thick, font=\small\bfseries},
  inetnode/.style={draw, ellipse, minimum width=24mm, minimum height=10mm, font=\scriptsize, fill=isLight, thick},
  process/.style={isnode, fill=isOT!15, draw=isOT, thick, minimum width=24mm},
  zone/.style={
    draw, rounded corners=4pt, thick, dashed, inner sep=4mm
  },
  conduit/.style={-{Stealth[length=2.5mm]}, thick, isNet!80!black},
  attack/.style={-{Stealth[length=2.5mm]}, thick, isAttack, dashed},
  flow/.style={-{Stealth[length=2.5mm]}, thick, isOT!70!black},
  netline/.style={thick, black!60},
  axisline/.style={-{Stealth[length=2mm]}, thick, black!70},
  tag/.style={font=\scriptsize\itshape, fill=white, inner sep=1pt},
  level/.style={
    draw, fill=isLight, minimum width=0.85\linewidth, minimum height=8mm,
    font=\small, anchor=west
  },
  brace/.style={decorate, decoration={brace, amplitude=4pt, raise=2pt}},
  legendbox/.style={draw, rounded corners=2pt, fill=isLight, inner sep=2mm, font=\scriptsize, align=left}
}

% --- handy macros ----------------------------------------------------
% \plcicon, \hmiicon etc as simple labelled boxes; useful inline.
\newcommand{\plcicon}[1][PLC]{\tikz[baseline=-0.6ex]{\node[plc, minimum width=10mm, minimum height=5mm, font=\scriptsize, inner sep=1pt]{#1};}}
\newcommand{\hmiicon}[1][HMI]{\tikz[baseline=-0.6ex]{\node[hmi, minimum width=10mm, minimum height=5mm, font=\scriptsize, inner sep=1pt]{#1};}}
\newcommand{\fwicon}[1][FW]{\tikz[baseline=-0.6ex]{\node[firewall, minimum width=10mm, minimum height=5mm, font=\scriptsize, inner sep=1pt]{#1};}}


\providecommand{\courseAuthor}{Industrial Security Lecture Series}
\providecommand{\courseInstitute}{Department of Computer Science}
\providecommand{\companyLogo}{logo}

\definecolor{slideAccent}{HTML}{00205B}
\definecolor{slideSecondary}{HTML}{00A0DC}

% --- Corporate branding overrides ------------------------------------
\IfFileExists{../../../branding.tex}{% =====================================================================
% Corporate / institutional branding -- single file to customise the
% look of the entire course (slides, notes, exercises, labs).
%
% HOW TO REBRAND IN UNDER A MINUTE:
%   1. Replace `branding/logo.pdf` with your own logo
%      (PDF preferred; PNG and JPG also work -- name it `logo.<ext>`).
%   2. (Optional) change the two colours below to your brand palette.
%   3. (Optional) override the author/institute lines below.
%   4. Re-run `make` at the repo root. That's it.
%
% Everything in this file has sensible defaults; you can delete a line
% to fall back to the built-in defaults, or comment it out.
%
% This file is loaded automatically by the slides, exercise, and lab
% preambles -- you never need to \input it manually.
% =====================================================================

% --- Logo --------------------------------------------------------------
% Path is resolved from each leaf-build directory; the default points at
% the shared `branding/` folder at the repo root. If you put your logo
% somewhere else, set the full or relative path here (without extension):
\renewcommand{\companyLogo}{../../../branding/logo}

% --- Author / institute (appears on the title page footer) ------------
\renewcommand{\courseAuthor}{}
\renewcommand{\courseInstitute}{}

% --- Brand colours ----------------------------------------------------
% slideAccent      = primary brand colour (title bands, accent rule)
% slideSecondary   = secondary brand colour (section badge, highlight)
% Both use HTML hex codes. Keep enough contrast against white text.
\definecolor{slideAccent}{HTML}{00205B}      % deep navy
\definecolor{slideSecondary}{HTML}{00A0DC}   % light blue accent
}{}

\colorlet{labAccent}{slideAccent}
\colorlet{labSec}{slideSecondary}
\definecolor{labCheck}{HTML}{2E7D32}
\definecolor{labWarn}{HTML}{B23A48}

\lstset{
  backgroundcolor=\color{black!4},
  basicstyle=\ttfamily\footnotesize,
  breaklines=true,
  frame=leftline,
  framerule=2pt,
  rulecolor=\color{labAccent},
  numbers=left,
  numberstyle=\tiny\color{black!50},
  showstringspaces=false,
  tabsize=2
}

\setlength{\tabcolsep}{4pt}

% Let TeX add a little extra inter-word stretch as a last resort before
% it reports an Overfull \hbox. Only kicks in on lines that would
% otherwise overflow (long \texttt paths, URLs, CIDRs); never loosens a
% line that already fits. Keeps prose/itemize within the margin.
\setlength{\emergencystretch}{3em}

\pagestyle{fancy}
\fancyhf{}
\renewcommand{\headrulewidth}{0.4pt}
\lhead{\small \courseTitle{} -- Lab \labNumber}
\rhead{\small Maps to Section \labSection}
\cfoot{\thepage}

\newtcolorbox{stepbox}[1]{
  enhanced, breakable, arc=2pt, boxrule=0pt,
  colback=labAccent!7, colframe=labAccent, leftrule=3pt,
  fonttitle=\bfseries\small, coltitle=white,
  colbacktitle=labAccent,
  title={\faTools\ Step #1}
}

\newtcolorbox{warnbox}{
  enhanced, breakable, arc=2pt, boxrule=0pt,
  colback=labWarn!7, colframe=labWarn, leftrule=4pt,
  fonttitle=\bfseries\small, coltitle=white, colbacktitle=labWarn,
  title={\faExclamationTriangle\ Caution}
}

\newtcolorbox{notebox}{
  enhanced, breakable, arc=2pt, boxrule=0pt,
  colback=labAccent!4, colframe=labAccent, leftrule=2pt,
  fonttitle=\bfseries\small, coltitle=white, colbacktitle=labAccent,
  title={\faInfoCircle\ Note}
}

\newtcolorbox{goalbox}{
  enhanced, breakable, arc=2pt, boxrule=0pt,
  colback=labCheck!7, colframe=labCheck, leftrule=3pt,
  fonttitle=\bfseries\small, coltitle=white, colbacktitle=labCheck,
  title={\faBullseye\ Goal}
}

\newtcolorbox{deliverbox}{
  enhanced, breakable, arc=2pt, boxrule=0pt,
  colback=labSec!10, colframe=labSec, leftrule=3pt,
  fonttitle=\bfseries\small, coltitle=white, colbacktitle=labSec,
  title={\faClipboardList\ Deliverable}
}

% --- Solution-sheet boxes (used only by lab-solutions.tex) ----------
\newtcolorbox{solutionbox}[1]{
  enhanced, breakable, arc=2pt, boxrule=0pt,
  colback=labCheck!7, colframe=labCheck, leftrule=3pt,
  fonttitle=\bfseries\small, coltitle=white, colbacktitle=labCheck,
  title={\faCheckCircle\ #1}
}

\newtcolorbox{rubricbox}{
  enhanced, breakable, arc=2pt, boxrule=0pt,
  colback=labSec!7, colframe=labSec, leftrule=3pt,
  fonttitle=\bfseries\small, coltitle=white, colbacktitle=labSec,
  title={\faBalanceScale\ Grading rubric}
}

\titleformat{\section}{\large\bfseries\color{labAccent}}{\thesection}{0.5em}{}

\newcommand{\labheader}{%
  \noindent
  \begin{tikzpicture}
    \fill[labAccent] (0,0) rectangle (\textwidth, -1.6cm);
    \node[anchor=north west, text=white, font=\bfseries\large] at (0.6cm, -0.4cm) {Lab \labNumber{} -- \labTitle};
    \node[anchor=south west, text=white, font=\small] at (0.6cm, -1.3cm) {\courseTitle{} \quad $\bullet$ \quad Maps to Section \labSection{} \quad $\bullet$ \quad \labTime};
    \node[anchor=north east, fill=labSec, text=white, font=\bfseries, inner sep=4pt, rounded corners=2pt]
      at ([xshift=-0.4cm, yshift=-0.4cm]\textwidth,0) {LAB \labNumber};
  \end{tikzpicture}
  \vspace{4mm}
}

% =====================================================================

\usepackage[utf8]{inputenc}
\usepackage[T1]{fontenc}
\usepackage[english]{babel}
\usepackage[a4paper, margin=2cm]{geometry}
\usepackage{graphicx}
\usepackage{listings}
\usepackage{xcolor}
\usepackage{colortbl}
\usepackage{tikz}
\usepackage{amsmath}
\usepackage{amssymb}
\usepackage{booktabs}
\usepackage{enumitem}
\usepackage{titlesec}
\usepackage{fancyhdr}
\usepackage{hyperref}
\usepackage{tcolorbox}
\tcbuselibrary{skins, breakable}
\usepackage{fontawesome5}

% =====================================================================
% Shared TikZ styles for the Industrial Security lecture series.
% Loaded by every slide deck and exercise sheet that uses TikZ.
% =====================================================================

\usetikzlibrary{
  arrows.meta,
  shapes,
  shapes.geometric,
  shapes.symbols,
  shapes.multipart,
  positioning,
  fit,
  calc,
  backgrounds,
  decorations.pathmorphing,
  decorations.pathreplacing,
  decorations.text,
  patterns,
  matrix,
  shadows
}

% --- colour palette --------------------------------------------------
\definecolor{isOT}{HTML}{1F4E79}     % OT (deep blue)
\definecolor{isIT}{HTML}{6F2DA8}     % IT (purple)
\definecolor{isSafety}{HTML}{C00000} % safety red
\definecolor{isNet}{HTML}{2E7D32}    % network green
\definecolor{isAttack}{HTML}{B23A48} % attack red
\definecolor{isAccent}{HTML}{E08E0B} % amber
\definecolor{isMute}{HTML}{6B6B6B}   % grey
\definecolor{isLight}{HTML}{EDF2F8}  % light background
\definecolor{isPLC}{HTML}{2C5282}
\definecolor{isHMI}{HTML}{2F855A}
\definecolor{isSCADA}{HTML}{6B46C1}

% --- generic node styles --------------------------------------------
\tikzset{
  isnode/.style={
    draw, rounded corners=2pt, minimum height=8mm, minimum width=18mm,
    font=\small, align=center, thick
  },
  isbox/.style={isnode, fill=isLight},
  plc/.style={isnode, fill=isPLC!15, draw=isPLC, thick},
  hmi/.style={isnode, fill=isHMI!15, draw=isHMI, thick},
  scada/.style={isnode, fill=isSCADA!15, draw=isSCADA, thick},
  sensor/.style={isnode, fill=isNet!10, draw=isNet, minimum height=6mm, minimum width=14mm, font=\scriptsize},
  actuator/.style={isnode, fill=isAccent!15, draw=isAccent, minimum height=6mm, minimum width=14mm, font=\scriptsize},
  firewall/.style={isnode, fill=isAttack!15, draw=isAttack, thick},
  server/.style={isnode, fill=isMute!15, draw=isMute!80, thick},
  switch/.style={isnode, fill=isLight, draw=black!60, minimum height=6mm, minimum width=14mm, font=\scriptsize},
  workstation/.style={isnode, fill=white, draw=isOT, thick},
  attacker/.style={isnode, fill=isAttack!20, draw=isAttack, thick, font=\small\bfseries},
  inetnode/.style={draw, ellipse, minimum width=24mm, minimum height=10mm, font=\scriptsize, fill=isLight, thick},
  process/.style={isnode, fill=isOT!15, draw=isOT, thick, minimum width=24mm},
  zone/.style={
    draw, rounded corners=4pt, thick, dashed, inner sep=4mm
  },
  conduit/.style={-{Stealth[length=2.5mm]}, thick, isNet!80!black},
  attack/.style={-{Stealth[length=2.5mm]}, thick, isAttack, dashed},
  flow/.style={-{Stealth[length=2.5mm]}, thick, isOT!70!black},
  netline/.style={thick, black!60},
  axisline/.style={-{Stealth[length=2mm]}, thick, black!70},
  tag/.style={font=\scriptsize\itshape, fill=white, inner sep=1pt},
  level/.style={
    draw, fill=isLight, minimum width=0.85\linewidth, minimum height=8mm,
    font=\small, anchor=west
  },
  brace/.style={decorate, decoration={brace, amplitude=4pt, raise=2pt}},
  legendbox/.style={draw, rounded corners=2pt, fill=isLight, inner sep=2mm, font=\scriptsize, align=left}
}

% --- handy macros ----------------------------------------------------
% \plcicon, \hmiicon etc as simple labelled boxes; useful inline.
\newcommand{\plcicon}[1][PLC]{\tikz[baseline=-0.6ex]{\node[plc, minimum width=10mm, minimum height=5mm, font=\scriptsize, inner sep=1pt]{#1};}}
\newcommand{\hmiicon}[1][HMI]{\tikz[baseline=-0.6ex]{\node[hmi, minimum width=10mm, minimum height=5mm, font=\scriptsize, inner sep=1pt]{#1};}}
\newcommand{\fwicon}[1][FW]{\tikz[baseline=-0.6ex]{\node[firewall, minimum width=10mm, minimum height=5mm, font=\scriptsize, inner sep=1pt]{#1};}}


\providecommand{\courseAuthor}{Industrial Security Lecture Series}
\providecommand{\courseInstitute}{Department of Computer Science}
\providecommand{\companyLogo}{logo}

\definecolor{slideAccent}{HTML}{00205B}
\definecolor{slideSecondary}{HTML}{00A0DC}

% --- Corporate branding overrides ------------------------------------
\IfFileExists{../../../branding.tex}{% =====================================================================
% Corporate / institutional branding -- single file to customise the
% look of the entire course (slides, notes, exercises, labs).
%
% HOW TO REBRAND IN UNDER A MINUTE:
%   1. Replace `branding/logo.pdf` with your own logo
%      (PDF preferred; PNG and JPG also work -- name it `logo.<ext>`).
%   2. (Optional) change the two colours below to your brand palette.
%   3. (Optional) override the author/institute lines below.
%   4. Re-run `make` at the repo root. That's it.
%
% Everything in this file has sensible defaults; you can delete a line
% to fall back to the built-in defaults, or comment it out.
%
% This file is loaded automatically by the slides, exercise, and lab
% preambles -- you never need to \input it manually.
% =====================================================================

% --- Logo --------------------------------------------------------------
% Path is resolved from each leaf-build directory; the default points at
% the shared `branding/` folder at the repo root. If you put your logo
% somewhere else, set the full or relative path here (without extension):
\renewcommand{\companyLogo}{../../../branding/logo}

% --- Author / institute (appears on the title page footer) ------------
\renewcommand{\courseAuthor}{}
\renewcommand{\courseInstitute}{}

% --- Brand colours ----------------------------------------------------
% slideAccent      = primary brand colour (title bands, accent rule)
% slideSecondary   = secondary brand colour (section badge, highlight)
% Both use HTML hex codes. Keep enough contrast against white text.
\definecolor{slideAccent}{HTML}{00205B}      % deep navy
\definecolor{slideSecondary}{HTML}{00A0DC}   % light blue accent
}{}

\colorlet{labAccent}{slideAccent}
\colorlet{labSec}{slideSecondary}
\definecolor{labCheck}{HTML}{2E7D32}
\definecolor{labWarn}{HTML}{B23A48}

\lstset{
  backgroundcolor=\color{black!4},
  basicstyle=\ttfamily\footnotesize,
  breaklines=true,
  frame=leftline,
  framerule=2pt,
  rulecolor=\color{labAccent},
  numbers=left,
  numberstyle=\tiny\color{black!50},
  showstringspaces=false,
  tabsize=2
}

\setlength{\tabcolsep}{4pt}

% Let TeX add a little extra inter-word stretch as a last resort before
% it reports an Overfull \hbox. Only kicks in on lines that would
% otherwise overflow (long \texttt paths, URLs, CIDRs); never loosens a
% line that already fits. Keeps prose/itemize within the margin.
\setlength{\emergencystretch}{3em}

\pagestyle{fancy}
\fancyhf{}
\renewcommand{\headrulewidth}{0.4pt}
\lhead{\small \courseTitle{} -- Lab \labNumber}
\rhead{\small Maps to Section \labSection}
\cfoot{\thepage}

\newtcolorbox{stepbox}[1]{
  enhanced, breakable, arc=2pt, boxrule=0pt,
  colback=labAccent!7, colframe=labAccent, leftrule=3pt,
  fonttitle=\bfseries\small, coltitle=white,
  colbacktitle=labAccent,
  title={\faTools\ Step #1}
}

\newtcolorbox{warnbox}{
  enhanced, breakable, arc=2pt, boxrule=0pt,
  colback=labWarn!7, colframe=labWarn, leftrule=4pt,
  fonttitle=\bfseries\small, coltitle=white, colbacktitle=labWarn,
  title={\faExclamationTriangle\ Caution}
}

\newtcolorbox{notebox}{
  enhanced, breakable, arc=2pt, boxrule=0pt,
  colback=labAccent!4, colframe=labAccent, leftrule=2pt,
  fonttitle=\bfseries\small, coltitle=white, colbacktitle=labAccent,
  title={\faInfoCircle\ Note}
}

\newtcolorbox{goalbox}{
  enhanced, breakable, arc=2pt, boxrule=0pt,
  colback=labCheck!7, colframe=labCheck, leftrule=3pt,
  fonttitle=\bfseries\small, coltitle=white, colbacktitle=labCheck,
  title={\faBullseye\ Goal}
}

\newtcolorbox{deliverbox}{
  enhanced, breakable, arc=2pt, boxrule=0pt,
  colback=labSec!10, colframe=labSec, leftrule=3pt,
  fonttitle=\bfseries\small, coltitle=white, colbacktitle=labSec,
  title={\faClipboardList\ Deliverable}
}

% --- Solution-sheet boxes (used only by lab-solutions.tex) ----------
\newtcolorbox{solutionbox}[1]{
  enhanced, breakable, arc=2pt, boxrule=0pt,
  colback=labCheck!7, colframe=labCheck, leftrule=3pt,
  fonttitle=\bfseries\small, coltitle=white, colbacktitle=labCheck,
  title={\faCheckCircle\ #1}
}

\newtcolorbox{rubricbox}{
  enhanced, breakable, arc=2pt, boxrule=0pt,
  colback=labSec!7, colframe=labSec, leftrule=3pt,
  fonttitle=\bfseries\small, coltitle=white, colbacktitle=labSec,
  title={\faBalanceScale\ Grading rubric}
}

\titleformat{\section}{\large\bfseries\color{labAccent}}{\thesection}{0.5em}{}

\newcommand{\labheader}{%
  \noindent
  \begin{tikzpicture}
    \fill[labAccent] (0,0) rectangle (\textwidth, -1.6cm);
    \node[anchor=north west, text=white, font=\bfseries\large] at (0.6cm, -0.4cm) {Lab \labNumber{} -- \labTitle};
    \node[anchor=south west, text=white, font=\small] at (0.6cm, -1.3cm) {\courseTitle{} \quad $\bullet$ \quad Maps to Section \labSection{} \quad $\bullet$ \quad \labTime};
    \node[anchor=north east, fill=labSec, text=white, font=\bfseries, inner sep=4pt, rounded corners=2pt]
      at ([xshift=-0.4cm, yshift=-0.4cm]\textwidth,0) {LAB \labNumber};
  \end{tikzpicture}
  \vspace{4mm}
}


\begin{document}
\labheader

\begin{goalbox}
Treat the lab Docker stack as a tiny plant. (a) Discover every host on the
\texttt{172.28.0.0/16} network and pin each one to a Purdue level.
(b) Probe which TCP ports are reachable from the \emph{student} container --
those are the \emph{conduits} that exist today.
(c) Capture one Modbus poll, count packets per direction, and decide whether
each source--destination pair is in policy.
(d) Write a deny-by-default \texttt{nftables} ruleset that enforces the zone
boundaries you drew.
\end{goalbox}

\begin{notebox}{}
\textbf{IT analogy first.} A Purdue level is just a \emph{DMZ tier}: L0--L2
is the inside LAN (databases, internal services), L3 is your application
tier, L3.5 is the classical DMZ (reverse proxies, jump hosts), L4--L5 is
the office / Internet edge. A \emph{zone} is an AWS VPC + security group; a
\emph{conduit} is one explicit security-group rule between two VPCs. The
only OT twist is that the wire itself is part of the boundary, so VLAN
hopping and switch-port misconfiguration become a security problem, not
just a networking problem.
\end{notebox}

\section*{Pre-flight}

\begin{stepbox}{0 -- Bring up the stack and shell into the student}
You did this in Lab 02; the same commands apply.
\begin{lstlisting}[language=bash]
cd bachelor/labs/_stack
docker compose up -d
docker compose ps          # every service should show "Up" / "running"
\end{lstlisting}

Every command below runs \emph{inside} the \texttt{student} container --
that container plays the role of an engineering workstation parked on the
plant network. Open a shell in it:

\begin{lstlisting}[language=bash]
docker compose exec student bash
# you are now 172.28.0.20 on the is-lab bridge network
ip -4 addr show eth0 | grep inet
ip route
\end{lstlisting}

Expected: one IPv4 address \texttt{172.28.0.20/16} on \texttt{eth0} and a
default route via \texttt{172.28.0.1} (the Docker bridge gateway). If you
do not see \texttt{172.28.0.20}, stop and fix the stack before continuing.
\end{stepbox}

\section*{Part A -- Host discovery and Purdue mapping}

\begin{stepbox}{1 -- Sweep the lab subnet with nmap}
From inside the student container:
\begin{lstlisting}[language=bash]
nmap -sn -n 172.28.0.0/24
\end{lstlisting}

\texttt{-sn} skips port scanning and just pings (ICMP + ARP on a local
bridge). \texttt{-n} disables DNS. You should see at least five hosts up:

\begin{lstlisting}
Nmap scan report for 172.28.0.1     # Docker bridge gateway
Nmap scan report for 172.28.0.5     # landing  (web overview)
Nmap scan report for 172.28.0.6     # openplc-init (sidecar, may be gone)
Nmap scan report for 172.28.0.10    # openplc  (soft PLC, Modbus)
Nmap scan report for 172.28.0.11    # opcua    (OPC UA server)
Nmap scan report for 172.28.0.12    # mqtt     (Mosquitto broker)
Nmap scan report for 172.28.0.20    # student  (you)
\end{lstlisting}

The \texttt{openplc-init} container is a short-lived bootstrap and may
already be \texttt{Exited}; nmap will then not list \texttt{.6}. That is
fine.
\end{stepbox}

\begin{stepbox}{2 -- Fingerprint each host (per-port banner)}
\texttt{-sV} probes the open ports and tries to identify the service. Run
it once per host so the output is readable:
\begin{lstlisting}[language=bash]
nmap -Pn -sV -p 1-10000 172.28.0.10   # openplc
nmap -Pn -sV -p 1-10000 172.28.0.11   # opcua
nmap -Pn -sV -p 1-10000 172.28.0.12   # mqtt
nmap -Pn -sV -p 1-10000 172.28.0.5    # landing
\end{lstlisting}

Roughly what you should see:

\begin{lstlisting}
172.28.0.10  502/tcp  open  modbus
             8080/tcp open  http-proxy   (OpenPLC web UI)
172.28.0.11  4840/tcp open  opcua-binary (open62541)
172.28.0.12  1883/tcp open  mqtt         (Mosquitto)
172.28.0.5   80/tcp   open  http         (nginx -- landing page)
\end{lstlisting}

\texttt{-Pn} skips host discovery (treat host as up); useful because some
containers may not respond to ICMP. If \texttt{-sV} guesses ``unknown''
just record the raw port -- the OPC UA fingerprint is sometimes blank.
\end{stepbox}

\begin{stepbox}{3 -- Map each host to a Purdue level}
On paper or in your notes, fill in the table below. Think about
\emph{function}, not IP.

\begin{center}
\renewcommand{\arraystretch}{1.25}
\begin{tabular}{|l|l|l|l|}
\hline
\textbf{Host} & \textbf{Role in the ``plant''} & \textbf{Purdue level} & \textbf{IT analogue} \\
\hline
172.28.0.10 (openplc)  & ?  & ?  & ? \\
172.28.0.11 (opcua)    & ?  & ?  & ? \\
172.28.0.12 (mqtt)     & ?  & ?  & ? \\
172.28.0.5  (landing)  & ?  & ?  & ? \\
172.28.0.20 (student)  & ?  & ?  & ? \\
\hline
\end{tabular}
\end{center}

\textbf{Hint.} OpenPLC controls a (simulated) physical process $\rightarrow$
that is L1 (basic control) and the simulated I/O behind it is L0. The OPC
UA server is the canonical SCADA/HMI plant-data broker $\rightarrow$ L2.
MQTT is normally an IIoT broker that aggregates plant data for upper tiers
$\rightarrow$ L3 or L3.5. The landing page is just an enterprise-facing
HTML overview $\rightarrow$ L4. The student container is an engineering
workstation; in a real plant it would sit at L2 or in the iDMZ.

\emph{One sentence per host explaining why.}
\end{stepbox}

\begin{stepbox}{4 -- Draw the zones-and-conduits diagram}
Take any tool you like -- pen and paper, \texttt{draw.io},
\texttt{excalidraw}, \texttt{mermaid} -- and draw the lab stack as zones
and conduits:

\begin{itemize}\setlength{\itemsep}{2pt}
  \item One rectangle per \emph{zone}, labelled with a Purdue level and an
        IEC 62443 SL-T target (1--4). At minimum: a \textbf{Control zone}
        (PLC + its I/O), a \textbf{Supervisory zone} (HMI / OPC UA), an
        \textbf{IIoT/iDMZ zone} (MQTT broker), an \textbf{Enterprise
        zone} (landing page).
  \item One arrow per \emph{conduit}. Label each arrow with the protocol
        and TCP port (e.g.\ ``Modbus/TCP 502'').
  \item Mark the conduit you would replace with a \textbf{hardware data
        diode} if this were a safety-critical plant.
\end{itemize}

\textbf{IT analogue:} this is the same diagram you would draw for a
three-tier web app -- web tier, app tier, DB tier -- except the boxes are
labelled by Purdue level and the only sanctioned arrows are conduits.
\end{stepbox}

\section*{Part B -- What is actually reachable?}

\begin{stepbox}{5 -- From the student, probe every service}
The student container can talk to \emph{everything} because the Docker
bridge does not enforce inter-container ACLs. Verify it:
\begin{lstlisting}[language=bash]
# TCP connect on every well-known industrial port
for h in 172.28.0.5 172.28.0.10 172.28.0.11 172.28.0.12; do
  echo "--- $h ---"
  nmap -Pn -n --open -p 80,502,1883,4840,8080 $h
done
\end{lstlisting}

You should see \texttt{.5:80}, \texttt{.10:502}, \texttt{.10:8080},
\texttt{.11:4840}, \texttt{.12:1883} all \textbf{open}. The student
container therefore has a path to \emph{every} Purdue level -- enterprise
web, IIoT broker, supervisory OPC UA, and L1 Modbus. \textbf{In a real
plant this is exactly what you want a firewall to prevent.}
\end{stepbox}

\begin{stepbox}{6 -- Enumerate today's conduits}
Each \texttt{open} port observed in step 5 is a conduit that exists right
now in the stack. Write them down in this form:

\begin{center}
\renewcommand{\arraystretch}{1.2}
\begin{tabular}{|l|l|l|l|l|}
\hline
\textbf{From zone} & \textbf{To zone} & \textbf{Protocol} & \textbf{Port} & \textbf{In policy?} \\
\hline
student (engineering WS) & openplc (L1) & Modbus/TCP    & 502  & ?  \\
student (engineering WS) & openplc (L1) & HTTP (web UI) & 8080 & ?  \\
student (engineering WS) & opcua  (L2)  & OPC UA binary & 4840 & ?  \\
student (engineering WS) & mqtt   (L3)  & MQTT          & 1883 & ?  \\
student (engineering WS) & landing(L4)  & HTTP          & 80   & ?  \\
\hline
\end{tabular}
\end{center}

\textbf{Question to answer in your notes:} of these five conduits, which
two are obviously legitimate for an engineering workstation, and which
three are red flags? (Hint: an engineer should never reach the
\emph{enterprise} web tier through the plant network, and they should not
have a TCP path to a public IIoT broker either.)
\end{stepbox}

\section*{Part C -- Capture a real conduit and count the packets}

\begin{stepbox}{7 -- Start tcpdump in one shell}
Open a \emph{second} shell into the student container (leave the first
one for the Modbus client):
\begin{lstlisting}[language=bash]
docker compose exec student bash
mkdir -p /labs/captures
tcpdump -i eth0 -nn -w /labs/captures/lab03.pcap host 172.28.0.10
\end{lstlisting}

\texttt{tcpdump} writes a pcap of every frame between this host and the
PLC. The capture is mounted to \texttt{bachelor/labs/\_stack/captures/}
on your host so you can open it in Wireshark later.
\end{stepbox}

\begin{stepbox}{8 -- Generate one Modbus exchange in the first shell}
Back in your first student shell:
\begin{lstlisting}[language=bash]
python3 /labs/scripts/modbus_read.py
\end{lstlisting}

Expected output: \texttt{Registers 0..9: [0, 0, 0, 0, 0, 0, 0, 0, 0, 0]}
(values depend on the running ladder).

Now stop tcpdump in the second shell with \texttt{Ctrl-C}. You should see
something like \texttt{14 packets captured} -- the TCP handshake (3) plus
the Modbus read request/response pair with their ACKs plus the FIN/ACK
teardown (4).
\end{stepbox}

\begin{stepbox}{9 -- Count packets per direction with tshark}
Still inside the student container:
\begin{lstlisting}[language=bash]
# total packets
tshark -r /labs/captures/lab03.pcap | wc -l

# packets student -> plc
tshark -r /labs/captures/lab03.pcap \
       -Y 'ip.src == 172.28.0.20 && ip.dst == 172.28.0.10' | wc -l

# packets plc -> student
tshark -r /labs/captures/lab03.pcap \
       -Y 'ip.src == 172.28.0.10 && ip.dst == 172.28.0.20' | wc -l

# only Modbus application-layer messages
tshark -r /labs/captures/lab03.pcap -Y 'mbtcp' \
       -T fields -e ip.src -e ip.dst -e mbtcp.trans_id -e modbus.func_code
\end{lstlisting}

Typical result: roughly equal counts in each direction (every TCP packet
is acknowledged), and exactly \emph{two} Modbus application-layer
messages -- function code \texttt{3} request (read holding registers) and
function code \texttt{3} response.

\textbf{Why this matters.} In a real plant, ``one read every 200\,ms''
becomes thousands of packets per minute on a single conduit. Knowing
\emph{request/response pairs vs.\ raw packet count} is the difference
between a useful baseline and a misleading one when you later write IDS
signatures (Lab 08).
\end{stepbox}

\begin{stepbox}{10 -- List unique source--destination pairs}
\begin{lstlisting}[language=bash]
tshark -r /labs/captures/lab03.pcap \
       -T fields -e ip.src -e ip.dst \
  | sort -u
\end{lstlisting}

Expected: exactly two lines -- \texttt{172.28.0.20 $\rightarrow$ 172.28.0.10}
and \texttt{172.28.0.10 $\rightarrow$ 172.28.0.20}. If you see more, you
either captured broadcasts or some other student talked to the PLC during
your capture window. Note any extra pair in your notes -- in a real
investigation, every unexpected pair is a potential lateral-movement
indicator.
\end{stepbox}

\section*{Part D -- Propose a firewall ruleset}

\begin{stepbox}{11 -- Draft a deny-by-default nftables ruleset}
\textbf{You do not apply this to the stack -- you write it.} The
exercise is to design the rules that \emph{would} enforce the zone
boundaries you drew in step 4.

\textbf{Policy in plain language:}
\begin{itemize}\setlength{\itemsep}{2pt}
  \item Allow the engineering workstation (\texttt{172.28.0.20}) to read
        Modbus on the PLC (\texttt{172.28.0.10:502/tcp}).
  \item Allow the engineering workstation to reach the OpenPLC engineering
        web UI on \texttt{172.28.0.10:8080/tcp}.
  \item Allow the supervisory zone (here also \texttt{172.28.0.20} -- a
        real plant would have a separate HMI host) to talk OPC UA to
        \texttt{172.28.0.11:4840/tcp}.
  \item Allow the PLC to push telemetry to the IIoT broker on
        \texttt{172.28.0.12:1883/tcp} -- \emph{outbound only from the
        PLC}, never inbound to it (iron rule: push from OT, pull from IT).
  \item Deny everything else, log the deny.
\end{itemize}

A minimal \texttt{nftables} sketch (write this in a file, do not run it):

\begin{lstlisting}[language=bash]
# /etc/nftables.conf  (proposed -- not applied)
table inet plant {
  chain forward {
    type filter hook forward priority 0; policy drop;

    ct state established,related accept

    # Engineering WS  -->  PLC (Modbus + engineering web)
    ip saddr 172.28.0.20 ip daddr 172.28.0.10 tcp dport { 502, 8080 } accept

    # Supervisory (same host in this lab)  -->  OPC UA server
    ip saddr 172.28.0.20 ip daddr 172.28.0.11 tcp dport 4840 accept

    # PLC pushes telemetry  -->  MQTT broker (one-way: PLC initiates)
    ip saddr 172.28.0.10 ip daddr 172.28.0.12 tcp dport 1883 \
        ct state new,established accept

    # Anything else: log and drop
    log prefix "plant-drop: " flags all counter drop
  }
}
\end{lstlisting}

\textbf{Things to notice} (write a one-line note for each):
\begin{itemize}\setlength{\itemsep}{2pt}
  \item The \emph{first} matching rule is \texttt{ct state established}.
        Without it, the response packet from the PLC back to the
        workstation is dropped and your TCP session never completes.
  \item There is no rule that allows the broker
        (\texttt{172.28.0.12}) to initiate anything inbound to the PLC.
        That is the ``push from OT, pull from IT'' iron rule from the
        slides expressed in iptables/nftables vocabulary.
  \item Every deny is logged. In a real plant that log is the input to
        the IDS in Lab 08.
\end{itemize}
\end{stepbox}

\begin{stepbox}{12 -- Same idea as an iptables ruleset (optional)}
If you grew up on iptables, here is the equivalent first three rules --
just for comparison. Do not apply.
\begin{lstlisting}[language=bash]
iptables -P FORWARD DROP
iptables -A FORWARD -m state --state ESTABLISHED,RELATED -j ACCEPT
iptables -A FORWARD -s 172.28.0.20 -d 172.28.0.10 -p tcp \
         -m multiport --dports 502,8080 -j ACCEPT
iptables -A FORWARD -s 172.28.0.20 -d 172.28.0.11 -p tcp \
         --dport 4840 -j ACCEPT
iptables -A FORWARD -s 172.28.0.10 -d 172.28.0.12 -p tcp \
         --dport 1883 -j ACCEPT
iptables -A FORWARD -j LOG --log-prefix "plant-drop: "
iptables -A FORWARD -j DROP
\end{lstlisting}
\end{stepbox}

\section*{Troubleshooting}

\begin{stepbox}{T -- What if it does not work}
\begin{itemize}\setlength{\itemsep}{2pt}
  \item \texttt{nmap -sn} returns nothing $\rightarrow$ you are running it
        from your host, not from inside the student container. Use
        \texttt{docker compose exec student bash} first.
  \item \texttt{nmap -sV} on \texttt{172.28.0.10:502} shows
        \texttt{filtered} $\rightarrow$ the openplc-init sidecar has not
        finished. Wait $\sim$20\,s after \texttt{up -d} or check
        \texttt{docker logs is-openplc-init}.
  \item \texttt{tcpdump: eth0: You don't have permission} $\rightarrow$
        the student container is missing \texttt{NET\_RAW}; re-run
        \texttt{docker compose up -d} so it picks up
        \texttt{cap\_add: NET\_RAW} from the compose file.
  \item \texttt{tshark} complains it cannot dissect
        \texttt{modbus/mbtcp} $\rightarrow$ on some builds the field
        is \texttt{mbtcp} not \texttt{modbus}; both filters work in the
        examples above.
  \item Your capture has zero packets $\rightarrow$ you ran
        \texttt{modbus\_read.py} \emph{before} tcpdump was listening.
        Order matters: start the dump first, then the client.
\end{itemize}
\end{stepbox}

\section*{Stretch goal}

\begin{stepbox}{S -- Policy audit from a fresh capture}
Generate a slightly noisier capture and audit every flow against your
own ruleset:
\begin{lstlisting}[language=bash]
# in one shell
tcpdump -i eth0 -nn -w /labs/captures/lab03-audit.pcap \
        'host 172.28.0.10 or host 172.28.0.11 or host 172.28.0.12'

# in a second shell, generate a mix
python3 /labs/scripts/modbus_read.py
python3 /labs/scripts/opcua_browse.py
python3 /labs/scripts/mqtt_publish.py  || true   # OK if broker silent

# stop the dump (Ctrl-C in shell 1), then:
tshark -r /labs/captures/lab03-audit.pcap \
       -T fields -e ip.src -e ip.dst -e tcp.dstport \
  | sort -u
\end{lstlisting}

For each unique \texttt{(src, dst, dport)} triple decide: would my
nftables rules from step 11 \textbf{accept} or \textbf{drop} this flow?
Anything that should be dropped but currently passes through the Docker
bridge is a real-world finding -- record it.
\end{stepbox}

\begin{deliverbox}
Hand in a single PDF with:
\begin{enumerate}\setlength{\itemsep}{2pt}
  \item The host table from step 3 (Purdue level + one-sentence
        justification per host).
  \item The zones-and-conduits diagram from step 4 (drawing or
        screenshot). Mark the data-diode candidate.
  \item The conduit table from step 6, with the ``in policy / red flag''
        column filled in.
  \item Your packet counts from step 9 (total, per-direction, Modbus
        application-layer) and the unique src/dst list from step 10.
  \item Your \texttt{nftables} ruleset from step 11, plus the three
        one-line notes at the bottom of the step.
\end{enumerate}
Optional: stretch-goal audit table from step S.
\end{deliverbox}

\section*{References}

\begin{itemize}\setlength{\itemsep}{2pt}
  \item T.J. Williams, \emph{The Purdue Enterprise Reference
        Architecture}, Computers in Industry, 1994.
  \item IEC 62443-3-2:2020, \emph{Security risk assessment for system
        design} (zones and conduits).
  \item IEC 62443-3-3:2013, FR\,5 -- \emph{Restricted Data Flow}.
  \item K. Stouffer et al., \emph{NIST SP 800-82 Rev.\,3 -- Guide to OT
        Security}, Sep 2023.
  \item Idaho National Lab, \emph{Improving Industrial Control System
        Cybersecurity with Defense-in-Depth Strategies}, 2016.
  \item Netfilter project, \texttt{nftables} wiki:
        \url{https://wiki.nftables.org/}.
\end{itemize}

\begin{warnbox}
Everything in this lab runs inside the lab containers. Do \emph{not}
point \texttt{nmap}, \texttt{tcpdump}, or any of these scripts at hosts,
PLCs, or networks that you do not own. On a production OT network, even
a passive \texttt{nmap -sV} can crash legacy equipment -- vendors have
documented PLCs that fault on a single TCP SYN to an unexpected port.
The Docker stack here is intentionally flat (no firewall between
containers) so you can \emph{see} the problem that real plants use the
techniques in this lab to fix.
\end{warnbox}
\end{document}
